\section{The Platonic Tradition}

We include Plato, Aristotle, Plotinus, and Thomas Aquinas as the pillars of the Platonic Tradition. This does not mean that they agree on every detail, but that they participate in a mutually comprehensible discussion. Our commitment is based on the following principles:

\begin{itemize}
\item \textbf{Metaphysical Positivism}. The spiritual world can be experienced in an analogous way to our experience of the sensual world. However, the former is not experienced as sensory images but rather as formless intuition. 
\item \textbf{Meditation}: Meditation is one stage of the process which is followed by contemplation and union. Unlike the idea held in the popular mind, for which meditation is a means for a better, more stress free personal life, we mean focused attention on a specific idea that leads to deeper understanding beyond what linear thinking can accomplish. 
\item \textbf{Ur-platonism}. There is a metaphysic that is common to different peoples in different traditions. In the west, this is represented by Platonism, properly understood. 
\end{itemize}
We discussed our experiences with our meditations on one of these topics:

\begin{itemize}
\item Developing self-awareness or remembering our first moment 
\item Re-experiencing Adam's moment of self-awareness 
\end{itemize}
Maximos the Confessor in his work\footnote{\url{https://gornahoor.net/?p=7568}} on the Kingdom of Heaven supplements this Tradition.

The fundamental conception is this, as Maximos mentions:

\begin{quotex}
The kingdom of heaven consists in possessing an inviolate and pre-eternal knowledge of created things through perceiving their inner essences as they exist in God. 

\end{quotex}
That is our starting point for our upcoming discussions on the various manifestations of Ur-Platonism. To wit, our understanding of the Kingdom of Heaven is founded on our understanding of essences, not on the material, sensual world.

There are obstacles to the meditations, that seemed difficult to articulate. It seems that we create a factitious mental world that obscures the ``real" world, which is described as ``beautiful". However, on deeper thought, the world, as it is manifested, is not really so beautiful and innocent. Perhaps, and this can be the subject for a new meditation, is that what is beautiful is actually essences. That is, if we can ``see" the essences, we would then experience beauty. The highest essence of all is the vision of the entire cosmos in the mind of God. The closer we can get to understanding that vision, the closer we are to Heaven.

The angelic and planetary stages filter this great idea as it becomes manifested. For example, different levels will extend the great cosmic idea into time, or into space, and so on. In the opposite direction, we can follow Maximos:

\begin{quotex}
It [the Kingdom of Heaven] is a state similar to that of the angels, attained by those who are saved.

\end{quotex}
In the opposite direction, we can ascend from the manifest through these higher states. Dante describes this journey through the purgatorical states, then through the planetary spheres, to the angels and ultimately to Heaven.

These topics came up about the meaning of the Eden story.

\begin{itemize}
\item Why was there one commandment for Adam? Why is a commandment necessary for a moral world of free agents? 
\item Why was it not good for Adam to be alone? Why is Eve necessary? Or more generally, why is there the other? 
\end{itemize}
As for the reason for the fall of Adam, can anyone develop his full potentials if he faces no challenges? What good, then, is immortality for its own sake? The Edenic state is too static, and ultimately, too little. Hence, a higher salvation is necessary in order to achieve Heaven. This led to the question of whether Christ would have come, even had there been no fall? The consensus of this moot point is ``yes", since salvation leads to the Kingdom of Heaven, a much more elevated state then Eden. Maximos emphasizes that the Kingdom is depends on the individual:

\begin{quotex}
the grace of the kingdom is given to all according to the quality and quantity of the righteousness that is in them.

\end{quotex}
The novel \emph{Zanoni} illustrates some of these ideas. Mejnour and Zanoni were initiates who discovered the secret of immortality. Mejnour is rather static. He tries to train a student, who fails the intricate and contradictory thought processes necessary to be initiated. We should be able to see this in our discussions: linear thinking does not seem to work. One idea, that might seem to be fruitful initially, leads to an impasse. Then something else becomes necessary.

Zanoni gives up immortality for Love, family, honor. He then enters the wheel of man, woman, life, death. But that seems superior to a life of immortality without purpose.

We see that there is not always a simple, convincing answer to the most important questions. Following one line of thought can lead to an impasse or even its opposite. When faced with such an opposition, always look for the neutralizing or equilibrating force. We learned that early on, but it seems to be forgotten. It is a fundamental notion in our way of thinking. The third force can be difficult to notice. Never forget to get back to the basics.

Recommended exercise: Meditation is both active and passive. Use the active imagination to place yourself into the situation, while simultaneously listening to what ideas may appear in consciousness.



\flrightit{Posted on 2020-10-06 by Admin }
